\dropcap{T}{he breath} of the kiln went out behind him and did not come back the same, the way a room keeps the shape of a man for a moment after he has left it. \person{Sella Vail} carried him out of the \place{Vail kiln-house} in a basket with a fold of wool over him, and the door did its old trick of not quite shutting, and the heat that had held him for a whole day went out into the yard as if it had never been meant to stay. Four days old, \person{Corin} thought, and already carried out of the only warm room he had known, with the place where the stone had been empty at his belly and the wool smelling of woodsmoke and of the woman who had laid it there. He could not lift his head to watch the roof go away. He listened to it instead. The kiln took one more breath behind the wall and missed the one after it, and the gap was the last thing of the house that came with him.

Being carried was its own instruction. He had been carried on stretchers enough times to read the hands that did it: the height they held the weight at, the way they took the corners, whether they watched the face or the wheels. Sella carried him neither like a sack nor like a child. She carried him like a pot that has been filled and must not be slopped, her forearm under the basket and her hand at the small of it, the wool keeping his cheek from the rail of her hip. He had once been a man who decided where the stretchers went. Now he was the weight, and the weight was not asked.

She had not touched the thing at her throat. He knew the habit well enough now: she put her hand to \lore{the widow's mark} before she opened the kiln and at no other time, and this morning the coals were only being banked and the door not opened, so the hand stayed where it was, on the handle of the basket. A lesser woman would have crossed herself and asked the fire to keep him. Sella Vail carried him down the step and did not look back at the house she had made warm for him.

\place{the Clay Yard} was the colour of the day, which was the colour of ashes shaken out of a grate. The sleet had gone off in the night and left a dry cold behind it, the kind that gets into the nose before it gets into the bones, and every breath he made under the wool came back at him as a small smoke. Oiled paper lifted and let go at the windows of the low houses, over and over, the way a man's thumb goes over the edge of a table while he waits for a number. The slip pits lay to his left with ice at their edges the colour of old bone. Beyond the last roofs \place{the Ilm} went under its skin of ice, and the whole district said, without saying it, that the ice was holding, which meant somebody had gone to look and come back with a number they did not give.

The cold got into him through the places the wool did not cover: his cheek where it had ridden back, the soft spot at his temple, the wrist lying outside the fold. He knew what cold did to a body before it killed it, the long surrender of the small vessels, the confusion that comes on like sleep, the afterdrop that takes a man down at the moment everybody thinks him warm. He had pulled men out of rivers in the old life and watched them get worse in a blanket. Here he was in the blanket and the getting worse was happening to him, and the only thing he could do about it was keep a count. The infant under the wool wanted his mother and wanted the breast and did not care about the small vessels at all.

The market came to him in layers, the way a room comes into focus for a man who has been sitting in the dark: first the smell, then the noise, then the things themselves. The smell of pigs came over the smell of clay, wet straw and warm dung and the sweet-sour of swill, and under that the smoke of barley bread from the bakehouses and, on the north wind, the heavy sweet smoke of tallow. Then the voices: the moot shouting at \place{the Cross}, a roar with edges on it, and under the roar the small separate sounds of a market setting itself up, a board dropped, a coin rung on a stone, a pig driven over wet boards with a weight behind it that he felt travel up through Sella's arm and into the basket. Then a bell, struck somewhere out of sight. It was not a church bell, though it went on longer than a church bell would and nobody in the yard crossed himself at it. He let it be a bell. There were things in this town that went under one name and did another work, and in four days he had learned to let them stand.

The bell went on at its own time, not with the voices and not against them, and nobody in the market seemed to need it to agree with anything. It was the first thing he heard that morning that had no work attached to it that he could see, and he listened to it longer than he needed to, the way a man listens to a sound in a ward because it is the only thing in the room that does not need him to answer it.

He counted because there was nothing else to count that belonged to him. \person{Daniel Marchetti} had counted the exits in every room he walked into, two in the bay, one behind the curtain, a fire door that stuck; it was a habit from a place where knowing the door mattered more than knowing the man. In the old life he had counted the beats of a pulse before he spoke to a resident, and the count had outlived the man and come to rest in a body that could not reach a cup. He counted the nails of a floor, the breaths of a fire through a wall, the steps of a woman carrying him, because these were his and nobody had asked for them and nobody could tax them. Sella's stride was long and even, the stride of a woman who had carried sacks of slip and knew what a weight would bear. He counted her steps out of the yard, up the short rise and onto the road that went down toward the water, and he came to a hundred and nine before the pace changed, and a hundred and nine was not four and four and a gap. It was not the kiln's breathing or the floor's nails or \person[Dessa Silt]{Dessa}'s nineteen out and nineteen back. It was a line, the same line her steps up the yard had been on the third morning, and it ended where the road began to go down. He put the hundred and nine in the private place beside the gap and left it there, unentered.

\scenebreak

Then the market gave him its own numbers to keep, and he kept them in the same place. Three pigs came past before they were weighed, then two more, then a man with a pole drove a lean one over the boards and it went clumsy on the wet wood and was hauled up by the ear. Five. Then six. A woman counting a heap of turnips said a figure out loud and another woman said a different figure and they went on with their hands in the turnips. He did not know what the figures bought or sold, and it did not matter; the numbers were everywhere and meant nothing he could spend, while the one number that was truly his was the one nobody had written down. He kept them side by side the way a man keeps the small change of a pocket and never counts it against the rent.

The woman at the stall had a slate propped against a crate with a row of figures on it in a hand he could not read. He tried, because he had once read a chart upside down across a bed and known what it said; the marks would not come into words for him, and the effort of it brought the small animal sound up in his throat. He gave the row its shape and let the shape go. There was writing all over this market and none of it had been written for him to read, and that, he was beginning to understand, was the ordinary condition of a person in a district.

The head went where it wanted. He had given up trying to hold it up on the second day and had learned only to time the small animal sound that came when it turned past the point where the neck would permit it, a little wet noise in the throat that was not a word and did not ask for anything, though the women always treated it as if it had. He had once been able to tell, from the sound a man made in the dark, where the pain sat and how long it had been there. Here the sound was his own and meant only that he was cold and awake and uncomfortable, and it got the wool adjusted and nothing else. In the basket the fold had ridden down off his face and the cold had got into his mouth, and a thread of drool went from his lip and would not be wiped. He tried to push the fold back with the hand that had once put a shoulder into its socket in the dark, and the hand was a small soft thing with milky nails that closed on air two inches from his face and came away. He had put a thousand shoulders back and he could not keep a blanket on his own chin. The joke of it had gone flat around the third day, but the exactness stayed, and the exactness was worse than the joke had ever been.

The fold came back over his face before he could be ashamed of it. Sella did it without looking, the way she had put her hand near the slate in \place{House Silt} without touching it, a movement that had already happened by the time he knew she had moved. She did not say his name and did not say any finished feeling over him. “Keep that off his face,” she said, to the wind, or to the wool, or to the boy; it was not clear which. “The wind's in his mouth. He'll take the cold in his throat.” The little-door burn closed on the fold and let it go, and the hand went back to the handle of the basket.

A woman at a tub looked up as they passed and went back to her work without appearing to have seen them. Another, straightening from a stall with a slate in her hand, said something about the boy that was a question. Sella did not answer the question. She answered the wool.

“The house brings him,” she said. “Keep his chin covered. That's the cold's way in.”

“Whose house, then?” the woman said.

“The house,” said Sella, and shifted the basket to her other arm.

“And the mother?”

“Account-day,” the woman said, as if that were an answer to something.

“Keep the fold over his ears,” Sella said. “He'll have the ache in them by noon.”

“Don't set him in the draught,” Sella added, and that was the end of it. The question lay where it had fallen, not answered and not asked again.

He had been put back under cloth in a public place by a hand that would not meet his eye and would not let him thank it. A kindness that cannot be answered is a kindness that cannot be refused; he had learned that under another roof, and here it was given again in the shape of a fold of wool, and there was no version of the world in which he could hand it back. The surgeon in him noted the pressure of her arm through the cloth, measured and unhurried, knowing what the thing would bear. The infant in him wanted the breast and was cold and wanted his mother, and the two wanted different things and both were louder than the man who was trying to keep the count.

She did not set the basket down where a slate lay open on a block beside the stall. She set it on her hip instead and stood a half-pace off from the slate, the way she had stood a half-pace off from the birth-slate in House Silt, and the boy felt the small careful distance and filed it with the other distances he had been made to keep: the nineteen steps, the end of the yard, the wall between a bed and a loft. A slate with writing open on it was a thing to stand off from. Nobody had told him why, and nobody had to.

\scenebreak

At \place{the yard well} the stone lip had a rim of ice on it, white in the morning and already chipped where somebody had broken water in the dark. A bucket stood on its side with a skin of ice in the bottom. The woman at the tub had her hands in water that steamed and then did not, and she rubbed the same patch of the same sheet over and over, which is what a woman does when she wants a morning to go by without being made to think about it. Barley smoke and tallow smoke lay in the low air together. A little further on someone had set a pot of glaze to warm and the lead in it gave off its sweet sickly smell, and the smell of the pigs came over that, and the clay came under everything. The oiled paper at the windows of House Silt lifted three doors up the yard, and the tile roof of \place{the smithy} held a crust of frost that had not yet been touched by the sun.

The market did not look at him. That was the worst of it, worse than the cold and worse than the noise: the square was filling with people and not one of them had yet found a reason to look into the basket. He had seen a room full of people not look at a man before, in a bay where the man on the table was past saving, and he had understood it then as a kind of professional mercy, the silence of people who had already done what they could. Here it was ordinary, and it was done to him. He had moved from a room where he was kept into a room where he would be counted, and in neither room could he say a word.

The shouting at the Cross went on resolving into something with a shape to it. The roar broke into separate voices and the separate voices into words he could not quite hold, and under them the bell went on being a bell and the pigs went on being driven over the boards. He was carried into the middle of it and put down at the edge of it, and the market closed around the basket the way a room closes around a man who has come in and cannot find the door he came by. He had spent his life reading rooms: where the senior stood, who would not meet an eye, which nurse had been crying in the supply room and tried to hide it with a clean apron. Rooms told you what they were before anybody spoke. This room told him nothing he could use. It was a room made for counting and for being counted, and he was the thing being carried through it, and the mistake of his old life, which he had made a hundred times and never noticed making, was to suppose that being the man who read the room meant he was not also one of the things in it. He had counted a hundred and nine steps to get here and five or six pigs and a gap that belonged to nobody, and not one of those numbers could be said out loud by a four-day-old in a basket under a fold of wool. \lore{The Hearth Drawing} was two days off, and \place{Low Kiln} said \lore{the Five} were three days off, and the room he was now in had already started counting him without asking his name.

The shouting went on. He was inside it. There was no part of it he could address and nobody in it who had asked him anything at all, and the new number stayed in his head where nobody had asked for it, a hundred and nine, his and untaxable, while the market made its ordinary morning noise around the basket and the fold of wool and the small hand that could not lift itself to push the wool back.
